\chapter{学习检查清单}

\section{入门阶段}

\begin{itemize}
  \item 能解释编译、链接和运行的区别。
  \item 能区分对象、值、名字、引用和指针。
  \item 能写出不越界的 \code{vector} 遍历。
  \item 能把读输入、处理数据、输出结果拆成不同函数。
  \item 能用 \code{const&} 避免只读场景的复制。
\end{itemize}

\section{中级阶段}

\begin{itemize}
  \item 能用构造函数建立类不变量。
  \item 能说明 RAII 如何处理提前返回和异常。
  \item 能根据需求选择 \code{vector}、\code{unordered_map}、\code{map} 等容器。
  \item 能把错误放进接口，而不是用魔法返回值。
  \item 能写出基本单元测试覆盖边界输入。
\end{itemize}

\section{高级阶段}

\begin{itemize}
  \item 能用 concept 约束模板参数。
  \item 能判断 ranges 视图是否存在悬空风险。
  \item 能用测量定位性能瓶颈，而不是凭直觉优化。
  \item 能用 mutex、atomic 和 condition variable 解释同步边界。
  \item 能拆分大型类，控制依赖方向。
\end{itemize}

\section{项目验收清单}

\begin{itemize}
  \item 每个公开函数都能说清楚所有权：观察、修改、保存副本还是接管资源。
  \item 每个错误路径都有测试或明确的人工验证步骤。
  \item 每个容器选择都能说出主要操作和复杂度理由。
  \item 每个类至少有一个围绕不变量的测试，而不只是 getter 测试。
  \item 每次性能优化都有优化前后的输入规模、构建模式和测量结果。
  \item 每个后台线程都有退出协议，程序关闭时不会依赖强制杀线程。
\end{itemize}

\section{读代码自检问题}

\begin{enumerate}
  \item 这段代码有没有隐藏的资源释放规则？
  \item 有没有返回局部对象的引用、保存短生命周期视图、或长期持有可能失效的迭代器？
  \item 错误是被吞掉、打印后继续，还是进入了接口？
  \item 测试是否覆盖空输入、单元素、边界值和失败路径？
  \item 如果数据规模扩大一百倍，主要成本会落在哪里？
\end{enumerate}
